\chapter{Introdução}
\label{chap:introducao}

A produção mundial de plásticos passou de cerca de 2 milhões de toneladas por ano em 1950 para 380 milhões de toneladas em 2015, e estima-se que 79\% de todos os resíduos plásticos já gerados tenham se acumulado em aterros ou no ambiente \cite{geyer2017}. Uma fração desse material chega ao mar: \citeonline{jambeck2015} estimaram que, apenas em 2010, entre 4,8 e 12,7 milhões de toneladas de resíduos plásticos gerados em áreas costeiras entraram no oceano. Uma vez no ambiente, o plástico se fragmenta em partículas cada vez menores sem se degradar por completo \cite{barnes2009,andrady2011}. As partículas menores que 5 mm, chamadas de microplásticos \cite{frias2019}, já foram encontradas em sedimentos, praias e na coluna d'água de todos os oceanos \cite{thompson2004,eriksen2014}.

A maior parte desse plástico não é descartada diretamente no mar. Ele é gerado em terra, onde as pessoas vivem, e é levado ao oceano principalmente pelos rios \cite{lebreton2017,schmidt2017,meijer2021}. Um rio integra tudo o que acontece na sua bacia hidrográfica: a chuva que cai sobre uma cidade escoa para os córregos, os córregos para o rio principal, e o rio desemboca na foz. Se o plástico que chega à foz vem das pessoas que vivem a montante, é natural esperar que a quantidade de plástico na foz cresça com a população da bacia. Essa expectativa está implícita nos principais modelos globais de emissão de plástico por rios, que usam a densidade populacional e a gestão de resíduos da bacia como variáveis explicativas \cite{lebreton2017,meijer2021}.

Testar essa relação diretamente, com medições de campo e não com modelos, exige responder a três perguntas que parecem simples mas não são. A primeira é \textbf{onde fica cada bacia}. As bacias precisam ser delineadas a partir do relevo, e isso exige um algoritmo que, dado um mapa de altitudes, atribua cada ponto de terra à saída no oceano para onde a água dali escoaria. A segunda é \textbf{quanta gente mora em cada bacia}, e em particular como tratar as cidades, que muitas vezes ficam sobre divisores de água e drenam para mais de uma bacia. Isso exige uma definição de cidade que não dependa de fronteiras administrativas, traçadas por motivos históricos e políticos e não pela forma da ocupação humana \cite{rozenfeld2008,rozenfeld2011}. A terceira é \textbf{quanto plástico chega à foz de cada bacia}. As medições de microplástico são pontuais, feitas em praias e estuários escolhidos por conveniência logística, e raramente coincidem com a foz de um rio.

\section{Objetivos}
\label{sec:objetivos}

O objetivo geral desta dissertação é testar se a população residente numa bacia hidrográfica está associada à concentração de microplástico medida na foz dessa bacia, e medir a forma dessa associação como uma lei de escala (alometria) na América do Norte e Central, dos Estados Unidos e sul do Canadá até o Canal do Panamá.

Os objetivos específicos são:

\begin{enumerate}
	\item delinear as bacias hidrográficas costeiras da região a partir do relevo global \gls{GEBCO}, usando o algoritmo baseado em percolação invasiva (\acrshort{IPBA}) com todas as saídas possíveis para o oceano;
	\item definir as manchas urbanas a partir da grade de população \gls{GPW} com o \gls{CCA}, calibrando os seus dois parâmetros pela lei de Zipf de forma independente da variável de resposta;
	\item calcular a população de cada bacia, repartindo entre as bacias as manchas urbanas que cruzam divisores de água;
	\item associar a cada bacia uma medição de microplástico na sua foz a partir do banco de microplásticos marinhos da \gls{NOAA}, avaliando de forma explícita como a distância aceita entre a foz e o ponto de coleta afeta o resultado;
	\item ajustar leis de escala entre população e microplástico por bacia, com intervalos de confiança e testes de significância que não dependam de suposições de normalidade.
\end{enumerate}

\section{Contribuições}
\label{sec:contribuicoes}

As contribuições deste trabalho são as seguintes.

\begin{itemize}
	\item \textbf{Um \textit{pipeline} reprodutível} que integra três bases públicas e independentes (relevo, população e microplástico) numa análise por bacia. Todo o processamento é feito em dois \textit{notebooks} executáveis, que baixam os dados e o código de terceiros em tempo de execução.
	\item \textbf{Delineamento de bacias sem lista manual de rios.} Todo pixel de terra vizinho ao oceano é tratado como saída (\textit{sink}), e as bacias de rio emergem naturalmente como as que acumulam mais área. As bacias principais são escolhidas por um critério de cobertura de área, sem limiar arbitrário de área mínima.
	\item \textbf{Calibração do CCA pela lei de Zipf com um algoritmo eficiente.} Uma varredura de 15 mil combinações de parâmetros, que exigiria 15 mil execuções do CCA, é resolvida com uma estrutura de união e busca que ativa os pixels em ordem decrescente de densidade. O resultado é idêntico ao do CCA em cada combinação e o custo é de algumas dezenas de execuções. O critério de escolha exige que o expoente fique próximo de 1 em vários tamanhos de amostra e impõe um piso físico à distância de ligação, o que evita soluções degeneradas.
	\item \textbf{Repartição de manchas urbanas entre bacias}, de modo que cada parte da cidade conte para a bacia para onde de fato drena.
	\item \textbf{Uma análise explícita da sensibilidade ao raio de aceitação} entre a foz e o ponto de coleta, que mostra quantas bacias ficam disponíveis em cada distância e como os expoentes e correlações variam com ela.
	\item \textbf{Uma avaliação honesta do resultado.} A dissertação mostra que a relação entre população e microplástico, embora positiva, é fraca e frágil com os dados atuais, e identifica as causas: a cobertura espacial das coletas, as contagens zero e a correlação entre população e área das bacias.
\end{itemize}

\section{Organização do texto}
\label{sec:organizacao}

O \autoref{chap:fundamentacao} apresenta os conceitos usados no trabalho: microplásticos e o seu transporte por rios, delineamento de bacias por percolação invasiva, a definição de cidades pelo CCA, a lei de Zipf, leis de potência e leis de escala urbana. O \autoref{chap:trabalhos-relacionados} discute os trabalhos que relacionam plástico e população, os que definem cidades a partir de dados em grade e os que aplicam leis de escala a fenômenos urbanos. O \autoref{chap:metodologia} descreve os dados, o recorte da região e cada etapa do \textit{pipeline}. O \autoref{chap:resultados} apresenta e discute os resultados, e o \autoref{chap:conclusao} resume as conclusões, as limitações e os trabalhos futuros. O \autoref{apendice:reprodutibilidade} descreve como reproduzir a análise, e o \autoref{apendice:bacias} lista as bacias usadas nos ajustes.
